% !TeX root = ../main.tex
\section{Valore assoluto e proprietà fondamentali}
% Fonte: Lezione 6, pp.~18--19

\begin{definizione}
Il \emph{valore assoluto}, o \emph{modulo}, di $x\in\R$ è
\[
|x|=\max\{x,-x\}=\begin{cases}
x,&x\ge0,\\
-x,&x<0.
\end{cases}
\]
\end{definizione}

La funzione $x\mapsto|x|$ ha dominio $\R$ e immagine $[0,+\infty)$,
è pari, decresce su $(-\infty,0]$ e cresce su $[0,+\infty)$.

\begin{schema}
\begin{tikzpicture}[x=1cm,y=.8cm]
  \draw[asse] (-3.2,0) -- (3.3,0) node[right] {$x$};
  \draw[asse] (0,-.3) -- (0,3.4) node[above] {$y$};
  \draw[intervallo] (-2.8,2.8) -- (0,0) -- (2.8,2.8);
  \node[coloreschema,above left] at (2.7,2.7) {$y=|x|$};
\end{tikzpicture}
\end{schema}

Per il modulo di una funzione la scelta del ramo dipende dal segno
di \emph{$f(x)$}, non dal segno di $x$:
\[
|f(x)|=\begin{cases}
f(x),&f(x)\ge0,\\
-f(x),&f(x)<0.
\end{cases}
\]
Graficamente si conserva la parte del grafico sopra l'asse $x$ e si
riflette rispetto a tale asse la parte che sta sotto.

\begin{proposizione}
Per ogni $a\in\R$ e ogni $b\ge0$ valgono
\begin{gather*}
a\le|a|,\qquad |a|=0\ \Longleftrightarrow\ a=0,\\
|a|\le b\ \Longleftrightarrow\ -b\le a\le b,\\
|a|\ge b\ \Longleftrightarrow\ a\ge b\ \text{ oppure }\ a\le-b.
\end{gather*}
Per ogni $a,b\in\R$ vale inoltre la \emph{disuguaglianza triangolare}
\[
|a+b|\le|a|+|b|.
\]
\end{proposizione}
\begin{proof}
Poiché $|a|=\max\{a,-a\}$, si ha sia $a\le|a|$ sia $-a\le|a|$;
il massimo è zero esattamente quando $a=0$.
La condizione $|a|\le b$ significa contemporaneamente $a\le b$ e
$-a\le b$, cioè $-b\le a\le b$. La condizione $|a|\ge b$ significa
che almeno uno tra $a$ e $-a$ è maggiore o uguale a $b$, dando la
seconda equivalenza.

Infine $-|a|\le a\le|a|$ e $-|b|\le b\le|b|$; sommando si ottiene
\[
-(|a|+|b|)\le a+b\le|a|+|b|.
\]
La prima equivalenza, con soglia non negativa $|a|+|b|$, conclude
la prova della disuguaglianza triangolare.
\end{proof}

Le versioni strette sono $|a|<b\iff-b<a<b$ e
$|a|>b\iff a>b\text{ oppure }a<-b$, per $b\ge0$.
Nelle disequazioni è essenziale distinguere l'intersezione delle due
condizioni nel primo caso dall'unione nel secondo.

\section{Distanza sulla retta e nozione di metrica}
% Fonte: Lezione 6, pp.~19--20

La distanza usuale fra due numeri reali è
\[
d(x,y)=|x-y|.
\]
È la lunghezza del segmento che li separa e non dipende dall'ordine
in cui si considerano i punti.

\begin{definizione}
Una \emph{distanza}, o \emph{metrica}, su un insieme $E$ (in
particolare su un sottoinsieme di $\R$) è una funzione
$d\colon E\times E\to[0,+\infty)$ tale che, per ogni $x,y,z\in E$,
\begin{enumerate}
  \item $d(x,y)\ge0$;
  \item $d(x,y)=0$ se e solo se $x=y$;
  \item $d(x,y)=d(y,x)$ \emph{(simmetria)};
  \item $d(x,y)\le d(x,z)+d(z,y)$ \emph{(disuguaglianza triangolare)}.
\end{enumerate}
\end{definizione}

\begin{proof}[Verifica per la distanza usuale]
Il modulo è non negativo e si annulla soltanto in zero: quindi
$|x-y|\ge0$ e $|x-y|=0\iff x=y$. Inoltre
$|x-y|=|-(y-x)|=|y-x|$, che dà la simmetria. Infine
\[
|x-y|=|(x-z)+(z-y)|\le|x-z|+|z-y|.
\]
Sono verificate tutte le proprietà della definizione.
\end{proof}
